\section{Robustness}\label{sec:robustness}

I organize robustness checks around three concerns:
measurement artifacts, sample dependence, and time
variation in the diagnostic coefficient.

\subsection{Measurement}

The Carlson--Parkin quantification depends on the
perception threshold $\delta$ and the treatment of
uncertain respondents. Table~\ref{tab:delta_sensitivity}
in Appendix~\ref{app:robustness} reports the diagnostic
coefficient under three alternative threshold values
($\delta \in \{0.3, 0.5, 0.7\}$). The sign is stable
across all values. The magnitude varies by less than
20\% from the baseline, confirming that the result is
not an artifact of the quantification method.

Replacing the Carlson--Parkin measure with the raw PBoC
expectation index produces the same qualitative pattern:
positive revisions predict negative subsequent forecast
errors (Table~\ref{tab:alt_measures}). Agreement across
two independently constructed measures indicates that
the diagnostic finding reflects household behavior
rather than a quantification artifact.

\subsection{Sample stability}

Dropping the COVID period (2020Q1--2020Q4) attenuates
the coefficient but preserves its sign
(Table~\ref{tab:subsample}). The result is not driven
by pandemic-era outliers, though significance weakens
in the smaller sample as expected.

A split-sample analysis divides the data at 2018Q1.
Both subsamples retain the negative diagnostic sign.
The post-2018 estimate is larger in absolute value,
suggestive of stronger overreaction in the later,
more volatile period. The smaller subsample sizes
($N \approx 16$ each) limit precision, so I treat
sign stability as the informative margin rather than
point estimates.

\subsection{Time-varying coefficient: Bayesian
  state-space model}

The Bayesian SSM (Appendix~\ref{app:ssm}) allows the
diagnostic coefficient to vary each quarter. The
coefficient $\beta_t$ is negative in 91\% of quarters
at the 90\% credible level. Prior sensitivity analysis
confirms that the sign and approximate magnitude are
stable across baseline, loose, and tight prior
specifications (Table~\ref{tab:prior_sensitivity}).

Figure~\ref{fig:beta_t_path} plots the posterior mean
and 90\% credible interval over the sample period. The
path becomes more negative during geopolitical episodes
(2016Q1, 2018Q3, 2020Q2) and relaxes toward zero in
calmer periods. This pattern is suggestive of
state-dependent overreaction, though the aggregate
sample is too short for sharp inference on interaction
effects.

\begin{figure}[htbp]
\centering
\includegraphics[width=0.84\textwidth]{beta_t.pdf}
\caption{Time-varying diagnostic intensity $\beta_t$:
  posterior mean and 90\% credible interval from the
  Bayesian state-space model.}
% duplicate label removed (stale file): fig:beta_t_path
\end{figure}

\subsection{CFPS measurement robustness}

Table~\ref{tab:cfps_robustness} reports the CFPS household
regression under five alternative specifications. The baseline
three-point coding ($\{-1,0,1\}$) yields $\hat{\beta} = -0.549$.
A binary rise indicator, a binary fall indicator, and dropping
the ``unchanged'' category all preserve the negative sign.
Column~(5) uses the full multi-wave panel ($N = 30{,}352$) with
province $\times$ wave fixed effects but no demographic controls,
confirming that the result does not depend on control availability
or a single survey wave. The coefficient ranges from $-0.52$ to
$-0.55$ across linear specifications, addressing the concern that
the finding is an artifact of a single coding convention applied
to a coarse qualitative variable.

\begin{table}[htbp]
\centering
\caption{CFPS Measurement Robustness}
\label{tab:cfps_robustness}
\small
\begin{tabular}{lccccc}
\toprule
 & (1) & (2) & (3) & (4) & (5) \\
 & Baseline & Binary rise & Binary fall & Drop unchanged & Full panel \\
\midrule
$\hat{\beta}$ & -0.5490*** & -1.4282*** & 0.9405*** & -0.5287*** & -0.5203*** \\
 & (0.0085) & (0.0221) & (0.0248) & (0.0065) & (0.0033) \\
\addlinespace
Observations & 13,031 & 13,031 & 13,031 & 6,649 & 30,352 \\
$R^2$ & 0.4559 & 0.4835 & 0.2463 & 0.7261 & 0.3861 \\
\addlinespace
Demographic controls & Yes & Yes & Yes & Yes & No \\
Province $\times$ wave FE & Yes & Yes & Yes & Yes & Yes \\
\bottomrule
\end{tabular}
\begin{tablenotes}
\small
\item \textit{Notes:} Dependent variable is the forecast error proxy (fe\_proxy). Column~(1) uses the baseline $\{-1,0,1\}$ revision coding. Column~(2) codes a binary indicator equal to one when a respondent revises upward to ``prices will rise.'' Column~(3) codes a binary indicator for downward revision to ``prices will fall.'' Column~(4) drops observations with no revision. Column~(5) uses the full multi-wave panel with province$\times$wave fixed effects but no demographic controls. Columns~(1)--(4) include province$\times$wave fixed effects and controls for age, education years, income, gender, and urban status. Column~(5) includes province$\times$wave fixed effects only. Standard errors (in parentheses) are clustered at the province level. $^{*}p<0.10$; $^{**}p<0.05$; $^{***}p<0.01$.
\end{tablenotes}
\end{table}

\subsection{Distributional evidence}

Table~\ref{tab:fe_by_revision} and
Figure~\ref{fig:fe_distribution} report the distribution of
forecast errors by revision category. Households that
revised upward have lower mean forecast errors than those
that revised downward. The difference is negative and
statistically significant. The distributional evidence
complements the regression slope by showing that the
revision-error relationship is not driven by outliers or
functional-form assumptions.

\begin{table}[htbp]
\centering
\begin{threeparttable}
\caption{Forecast-Error Distribution by Revision Category}
\label{tab:fe_by_revision}
\begin{tabular}{lcccc}
\toprule
Category & $N$ & Mean & Median & SD \\
\midrule
Downward & 7,638 & 0.7426 & 1.0000 & 0.4372 \\
Unchanged & 14,983 & 0.4334 & 1.0000 & 0.7560 \\
Upward & 7,731 & -0.5645 & -1.0000 & 0.4959 \\
\midrule
Difference (Upward $-$ Downward) & & -1.3071*** & & \\
\bottomrule
\end{tabular}
\begin{tablenotes}[flushleft]\small
\item \textit{Notes:} The table reports summary statistics of the forecast-error proxy by belief-revision category. Downward: revision $< 0$; Unchanged: revision $= 0$; Upward: revision $> 0$. The difference row reports the mean difference with significance from a Welch two-sample $t$-test. $^{*}p<0.10$, $^{**}p<0.05$, $^{***}p<0.01$.
\end{tablenotes}
\end{threeparttable}
\end{table}


\begin{figure}[htbp]
\centering
\includegraphics[width=0.84\textwidth]{fe_distribution_by_revision.pdf}
\caption{Forecast-error density by revision category.
  Upward revisers (red) have lower subsequent forecast
  errors than downward revisers (blue).}
\label{fig:fe_distribution}
\end{figure}

Additional robustness checks, including BVAR structural
consistency, replication parity, control sensitivity,
and IV diagnostics, are reported in the online appendix.


